\chapter{Architectural Decision Framework and Discussion}
\label{chap:discussion}

\section{Hypothesis Evaluation}
\label{sec:hypotheses-eval}

Evaluating the empirical evidence from Chapter~4 against the three research hypotheses formulated in the Introduction yields the following outcomes:
\begin{itemize}[leftmargin=1.5em]
  \item \textbf{H1 (Confirmed):} The modular monolith consumed $25.5\%$ less memory at idle ($903.98\text{ MiB}$ vs.\ $1,213.60\text{ MiB}$) and $14.7\%$ less memory at $20\text{ VUs}$ ($1,151.35\text{ MiB}$ vs.\ $1,350.20\text{ MiB}$), confirming that process consolidation yields substantial memory savings on SME hardware.
  \item \textbf{H2 (Confirmed):} While aggregate throughput was comparable ($7.49\text{ req/s}$ vs.\ $7.58\text{ req/s}$), microservices achieved $2.65\times$ lower $p_{95}$ tail latency ($12.30\text{ ms}$ vs.\ $32.57\text{ ms}$), validating the responsiveness advantage of decoupled read endpoints.
  \item \textbf{H3 (Confirmed):} Single-domain scenarios (S1, S2, S4) touched fewer files per service in microservices ($2$--$3$ files) with single-container redeployment, whereas the cross-cutting multi-currency scenario (S5) required $9$ file edits across $3$ deployment units in microservices versus $6$ files in a single deployment unit in the monolith.
\end{itemize}

\section{SME Architectural Decision Matrix}
\label{sec:decision-matrix}

Combining quantitative metrics with practitioner insights produces a structured decision matrix (Table~\ref{tab:decision-matrix}) for SME technical leaders.

\begin{table}[h]
  \caption{Recommended architectural selection matrix for SME e-commerce enterprises}
  \label{tab:decision-matrix}
  \centering
  \begin{tabular}{p{3.4cm}p{5.2cm}p{5.2cm}}
    \toprule
    \textbf{Decision Factor} & \textbf{Choose Modular Monolith} & \textbf{Choose Selective Microservices} \\
    \midrule
    Team Structure & 1--6 engineers sharing a unified backlog & $\ge 2$ autonomous squads owning distinct bounded contexts \\
    Hosting Budget & Single VPS / container node ($\le 4\text{ GiB}$ RAM) & Managed orchestrator with dedicated observability stack \\
    Domain Stability & Rapidly evolving business rules and cross-cutting features & Mature, stable domain contracts with isolated scaling hotspots \\
    Latency SLA & Standard storefront responsiveness ($p_{95} < 50\text{ ms}$) & Strict sub-$15\text{ ms}$ API SLAs under burst concurrency \\
    Recommended Path & \textbf{Default starting architecture} with strict DDD module boundaries & \textbf{Incremental extraction} (Strangler Fig) of Catalog/Search first \\
    \bottomrule
  \end{tabular}
\end{table}

\section{Threats to Validity and Limitations}
\label{sec:threats}
\begin{itemize}[leftmargin=1.5em]
  \item \textbf{Internal Validity:} All containers ran on a single local WSL2 host rather than multi-node cloud VMs; while this eliminated external network variance, inter-container loopback latency in production clusters would add $1$--$3\text{ ms}$ per network hop to microservices.
  \item \textbf{External Validity:} Findings are grounded in the .NET / ASP.NET Core ecosystem; interpreted runtimes (e.g., Node.js, Python) or ahead-of-time native binaries (Go, Rust) will exhibit different absolute memory baselines while preserving the relative architectural trade-offs.
  \item \textbf{Construct Validity:} Change propagation was measured via modified source files and deployment units; incorporating longitudinal defect rates across multi-month sprints remains an avenue for future work.
\end{itemize}

\section{Chapter Summary}
\label{sec:ch5-summary}
This chapter confirmed Hypotheses H1--H3, synthesized the findings into a practical five-factor SME architectural decision matrix favoring a ``Modular-Monolith-First'' trajectory, and documented internal, external, and construct validity boundaries.
